%% Drafted from the mapped corpus by scripts/draft_topics.py.
% Model: gpt-5.6-terra; source characters: 24927.
\begin{konspekt}
\kselect{На изпита ще се падне \textbf{или} алгоритъмът на Дийкстра \emph{и} алгоритъмът за тегловен ацикличен граф, \textbf{или} алгоритъмът на Белман--Форд \emph{и} алгоритъмът на Флойд--Уоршал. Двойките са неразделни.}
\kreq{Определение и варианти на задачата}{плюс потенциалните проблеми при отрицателни тегла --- \kref{8.1}.}
\kreq{Вариант, псевдокод, коректност, сложност при различни структури}{алгоритъм на Дийкстра --- \kref{8.3}.}
\kreq{Вариант, псевдокод, коректност, сложност, предимства пред Дийкстра}{най-къси пътища в тегловен ориентиран ацикличен граф --- \kref{8.4}.}
\kreq{Вариант, псевдокод, коректност, сложност, откриване на отрицателен цикъл}{алгоритъм на Белман--Форд --- \kref{8.5}.}
\kreq{Вариант, псевдокод, коректност, сложност по време и по памет}{алгоритъм на Флойд--Уоршал --- \kref{8.6}.}
\kreq{Литература}{[31].}
\end{konspekt}

\section{Постановка на задачата}\label{s:8.1}

\begin{defn}[Тегловен граф и тегло на маршрут]
Нека $G=(V,E)$ е тегловен граф с теглова функция
\[
\omega:E\to\mathbb{R}.
\]
Разглеждаме преди всичко ориентирани графи. Теглото на маршрут $p=(v_0,e_1,v_1,\ldots,e_k,v_k)$ е сумата от теглата на ребрата му, броени с повторенията:
\[
\omega(p)=\sum_{i=1}^{k}\omega(e_i).
\]
\end{defn}

\begin{defn}
Нека $u,v\in V$. С $\operatorname{Paths}(u,v)$ означаваме множеството от пътища от $u$ до $v$. Най-късото разстояние от $u$ до $v$ се означава с $\delta(u,v)$ и е
\[
\delta(u,v)=
\begin{cases}
+\infty, & \text{ако }\operatorname{Paths}(u,v)=\emptyset,\\
-\infty, & \text{ако е изпълнено }C(u,v),\\
\min\{\omega(p)\mid p\in\operatorname{Paths}(u,v)\}, & \text{иначе}.
\end{cases}
\]
Тук $C(u,v)$ означава, че има маршрут от $u$ до $v$, минаващ през отрицателен цикъл. Допускат се повторения на върхове и ребра; при липса на такъв цикъл минимумът се достига от прост път. Под \emph{най-къс път} ще разбираме път с най-малко тегло.
\end{defn}

\begin{defn}
В ориентиран тегловен граф \emph{отрицателен цикъл} е цикъл, чиято сума от теглата на ребрата е отрицателна.
\end{defn}

\begin{supp}[Бележка]
При неориентиран граф всяко ребро с отрицателно тегло позволява обхождане в двете посоки с отрицателно общо тегло. Следователно наличието на отрицателно ребро е съществено препятствие за стандартното понятие за най-къс път, дори ако се разглеждат само прости цикли.
\end{supp}

\subsection{Варианти на задачата}\label{s:8.1.1}

Основните варианти на задачата за най-къси пътища са следните.

\begin{itemize}
\item \emph{Най-къс път между дадена двойка върхове} (\emph{single-pair shortest paths}). Дадени са $G$, $\omega$, начален връх $s$ и краен връх $t$. Търси се $\delta(s,t)$.

\item \emph{Най-къси пътища от един източник} (\emph{single-source shortest paths}, SSSP). Дадени са $G$, $\omega$ и начален връх $s$. Търсят се всички стойности
\[
\delta(s,v),\qquad v\in V.
\]

\item \emph{Най-къси пътища до един краен връх} (\emph{single-destination shortest paths}). Дадени са $G$, $\omega$ и краен връх $t$. Търсят се всички стойности
\[
\delta(v,t),\qquad v\in V.
\]

\item \emph{Най-къси пътища между всички двойки върхове} (\emph{all-pairs shortest paths}, APSP). Дадени са $G$ и $\omega$. Търсят се всички стойности
\[
\delta(u,v),\qquad u,v\in V.
\]
\end{itemize}

Задачата с един краен връх може да се свежда до задача с един източник чрез обръщане на всички ребра. Задачата за една двойка е частен случай на задачата с един източник, но специализиран алгоритъм може да прекрати работата си по-рано, когато крайният връх е достатъчно обработен.

\subsection{Отрицателни тегла и оптимална подструктура}\label{s:8.1.2}

Отрицателните тегла на ребрата сами по себе си не правят задачата безсмислена. Проблемът възниква при отрицателен цикъл, който е достижим по път от началния до крайния връх. Тогава цикълът може да се обходи произволно много пъти и да се получават пътища с все по-малко тегло. Поради това няма път с минимално тегло и стойността е $-\infty$.

\begin{remark}
Ако от началния връх не е достижим отрицателен цикъл, за всеки достижим краен връх може да се избере прост най-къс път. Наистина, ако най-къс маршрут съдържа цикъл с неотрицателно тегло, премахването му не увеличава теглото; отрицателен цикъл по маршрута би бил достижим от началния връх. При цикъл с нулево тегло може да има и непрости най-къси маршрути, но винаги съществува прост представител със същото тегло.
\end{remark}

\begin{thm}[Оптимална подструктура на най-късите пътища]
Нека $p$ е най-къс път от $s$ до $t$ и нека $u$ предхожда $v$ по $p$. Тогава подпътят на $p$ от $u$ до $v$ е най-къс път от $u$ до $v$.
\end{thm}

\begin{proof}
Нека подпътят от $u$ до $v$ в $p$ е $q$. Ако съществува път $q'$ от $u$ до $v$ с
\[
\omega(q')<\omega(q),
\]
то можем да заменим $q$ с $q'$ в пътя $p$. Полученият път от $s$ до $t$ има тегло
\[
\omega(p)-\omega(q)+\omega(q')<\omega(p),
\]
което противоречи на избора на $p$.
\end{proof}

\section{Релаксация на ребра}\label{s:8.2}

Алгоритмите за най-къси пътища поддържат за всеки връх $v$ две полета:

\begin{itemize}
\item оценка $v.d$ за разстоянието от началния връх $s$ до $v$; тя е тегло на някакъв вече намерен маршрут от $s$ до $v$, тоест горна граница за $\delta(s,v)$;
\item предшественик $v.\pi$ --- върхът, от който този маршрут влиза в $v$. Пътят се възстановява, като се тръгне от $v$ и се следват указателите $\pi$ назад до $s$.
\end{itemize}

Използват се още две означения: $w((x,y))$ е теглото на реброто $(x,y)$, а $\mathit{adj}(x)$ е списъкът на съседство на $x$, тоест списъкът на върховете $y$, за които $(x,y)\in E$.

Инициализацията е следната; името \verb|ISS| идва от \emph{initialize single source}.

\begin{verbatim}
ISS(G, s)
    for each v in V(G)
        v.d  <- infinity       // още не е намерен маршрут до v
        v.pi <- Nil            // предшественик по този маршрут
    s.d <- 0                   // до източника води празният маршрут
\end{verbatim}

Основната операция е \emph{релаксация}.

\begin{verbatim}
Relax((x, y))
    if y.d > x.d + w((x, y))   // през x се стига по-евтино до y
        y.d  <- x.d + w((x, y))  // оценката се подобрява
        y.pi <- x                // и се запомня как е постигната
\end{verbatim}

Тя проверява дали вече известният път до $x$, последван от реброто $(x,y)$, подобрява текущата оценка за $y$. Релаксацията никога не увеличава оценка и никога не поставя стойност, която не е тегло на реален маршрут; сравнението със сума, съдържаща $x.d=+\infty$, е винаги лъжа, така че недостижимите досега върхове не се променят.

Всички алгоритми по-долу се различават само по това \emph{в какъв ред} релаксират ребрата; самата операция е една и съща.

\begin{lem}[Неравенство за ребро]
За всяко ребро $(x,y)\in E$ е изпълнено
\[
\delta(s,y)\leq \delta(s,x)+\omega((x,y)).
\]
\end{lem}

\begin{proof}
Ако $x$ не е достижим от $s$, дясната страна е $+\infty$ и неравенството е вярно. Ако до $x$ може да се достигне чрез път, съдържащ отрицателен цикъл, тогава и до $y$ може да се достигне по такъв път и и двете релевантни стойности са $-\infty$.

В останалия случай съществува най-къс път от $s$ до $x$. Добавянето на реброто $(x,y)$ дава път от $s$ до $y$ с тегло
\[
\delta(s,x)+\omega((x,y)).
\]
По определение $\delta(s,y)$ не може да е по-голямо от теглото на този път.
\end{proof}

\begin{lem}[Свойство на оценките]
След инициализация и след всяка произволна последователност от релаксации за всеки връх $v$ е изпълнено
\[
v.d\geq \delta(s,v).
\]
Освен това, щом за даден връх се достигне равенството
\[
v.d=\delta(s,v),
\]
то то се запазва при всички следващи релаксации.
\end{lem}

\begin{proof}
Твърдението се доказва с индукция по броя извършени релаксации. След инициализацията $s.d=0$, а за останалите върхове оценката е $+\infty$, следователно твърдението е вярно.

При релаксация на $(x,y)$ може да се промени само $y.d$. Ако се промени, новата му стойност е
\[
x.d+\omega((x,y))
 \geq \delta(s,x)+\omega((x,y))
 \geq \delta(s,y),
\]
където последното неравенство е от предишната лема. Релаксацията никога не увеличава оценки. Следователно достигнато точно разстояние не може по-късно да бъде развалено.
\end{proof}

\begin{lem}[Релаксация по най-къс път]
Нека $p$ е най-къс път от $s$ до $v$ и нека $(u,v)$ е последното му ребро. Ако преди релаксацията на $(u,v)$ е изпълнено
\[
u.d=\delta(s,u),
\]
то след тази релаксация е изпълнено
\[
v.d=\delta(s,v).
\]
\end{lem}

\begin{proof}
По оптималната подструктура префиксът на $p$ от $s$ до $u$ е най-къс път, следователно
\[
\delta(s,v)=\delta(s,u)+\omega((u,v)).
\]
След релаксацията
\[
v.d\leq u.d+\omega((u,v))=\delta(s,v).
\]
От свойството на оценките имаме и $v.d\geq\delta(s,v)$, откъдето следва равенство.
\end{proof}

\section{Алгоритъм на Дийкстра}\label{s:8.3}

\subsection{Вариант на задачата и идея}\label{s:8.3.1}

Алгоритъмът на Дийкстра решава задачата SSSP в ориентиран тегловен граф, когато всички тегла са неотрицателни:
\[
\omega(e)\geq 0,\qquad e\in E.
\]
Той поддържа множество $D$ от окончателно обработени върхове. На всяка стъпка се избира необработен връх с най-малка текуща оценка и той се добавя в $D$. След това се релаксират всички изходящи от него ребра.

\subsection{Използвани операции}\label{s:8.3.2}

Псевдокодът извиква операции на приоритетна опашка. Те не са част от самия алгоритъм: всяка реализация, която ги поддържа, е допустима, а изборът на реализация влияе само на сложността.

\begin{defn}[Приоритетна опашка по минимум]
Приоритетна опашка е структура от данни, която пази множество от елементи, всеки с числов ключ, и поддържа операциите:

\begin{itemize}
\item \verb|Insert(Q,x)| --- добавя елемента $x$ с текущия му ключ;
\item \verb|Extract-Min(Q)| --- премахва от $Q$ и връща елемент с минимален ключ;
\item \verb|Decrease-Key(Q,x,k)| --- заменя ключа на $x$ с новата стойност $k$, която не е по-голяма от старата, и възстановява вътрешния ред на структурата.
\end{itemize}
\end{defn}

В алгоритъма на Дийкстра ключът на връх $v$ е текущата оценка $v.d$. Записът $Q\leftarrow V(G)$ означава построяване на опашка от всички върхове с текущите им ключове, тоест $|V|$ операции \verb|Insert| или еднократно построяване.

\subsection{Псевдокод}\label{s:8.3.3}

\begin{verbatim}
Dijkstra(G, w, s)
    ISS(G, s)                  // всички оценки са infinity, s.d = 0
    D <- emptyset              // окончателно обработените върхове
    create priority queue Q    // ключ на връх v е оценката v.d
    Q <- V(G)                  // всички върхове чакат в опашката
    while Q is not empty
        x <- Extract-Min(Q)    // необработен връх с най-малък d
        D <- D union {x}       // от този момент x.d е окончателно
        for each y in adj(x)   // релаксация на изходящите ребра
            if y in Q and y.d > x.d + w((x, y))
                y.d  <- x.d + w((x, y))
                y.pi <- x
                Decrease-Key(Q, y, y.d)   // синхронизира опашката
\end{verbatim}

Означенията в псевдокода са следните.

\begin{itemize}
\item $G$, $w$ и $s$ са входът: ориентиран граф, теглова функция с неотрицателни стойности и начален връх.
\item $Q$ съдържа необработените върхове, а $D=V\setminus Q$ --- обработените; двете множества винаги се допълват, така че $D$ се пази само за удобство на изложението и може да се пропусне.
\item Тялото на \verb|while| се изпълнява по веднъж за всеки връх: \verb|Extract-Min| премахва връх и нищо не се връща обратно в $Q$.
\item Трите реда в тялото на \verb|if| са точно \verb|Relax((x,y))|, следвана от \verb|Decrease-Key|; последната само съобщава на опашката, че ключът на $y$ е намалял.
\end{itemize}

Проверката $y\in Q$ не позволява промяна на вече окончателно обработен връх; тя се реализира за константно време с отделна булева информация дали върхът е в опашката. Резултатът от алгоритъма са оценките $v.d$ и указателите $v.\pi$, които образуват дърво на най-късите пътища с корен $s$.

\subsection{Пример}\label{s:8.3.4}

\begin{example}
Разглеждаме графа от фигура~\ref{fig:dijkstra-example} с начален връх $s$.

\begin{lecfigure}[Ориентиран граф с неотрицателни тегла. Удебелените ребра са ребрата $\{(v.\pi,v)\}$ на полученото дърво на най-късите пътища.\label{fig:dijkstra-example}]
\begin{tikzpicture}[scale=1.25,
  vtx/.style={circle, draw=defframe, fill=defback, inner sep=0pt,
              minimum size=6.5mm, font=\small},
  earc/.style={gray!65, line width=0.5pt, -{Stealth[length=2mm]}},
  etree/.style={thmframe, line width=1.5pt, -{Stealth[length=2.4mm]}},
  wlab/.style={font=\footnotesize, inner sep=1.4pt, fill=white}]
  \node[vtx] (s) at (0,1)     {$s$};
  \node[vtx] (a) at (1.9,2.1) {$a$};
  \node[vtx] (b) at (1.9,-0.1){$b$};
  \node[vtx] (c) at (3.9,1.6) {$c$};
  \node[vtx] (d) at (5.8,0.7) {$d$};
  \draw[earc]  (s) -- node[wlab] {4}  (a);
  \draw[etree] (s) -- node[wlab] {2}  (b);
  \draw[etree] (b) -- node[wlab] {1}  (a);
  \draw[etree] (a) -- node[wlab] {5}  (c);
  \draw[earc]  (b) -- node[wlab] {8}  (c);
  \draw[earc]  (b) -- node[wlab] {10} (d);
  \draw[etree] (c) -- node[wlab] {2}  (d);
\end{tikzpicture}
\end{lecfigure}

След \verb|ISS| е $s.d=0$, а всички останали оценки са $\infty$. Всеки ред от таблицата отговаря на една итерация на цикъла \verb|while|: извлеченият връх, окончателната му оценка и състоянието на останалите оценки след релаксациите. Записът $3(b)$ означава $d=3$ и $\pi=b$.

\begin{center}
\begin{tabular}{llll}
\toprule
Извлечен $x$ & $x.d$ & Релаксирани ребра & Оценки в $Q$ след стъпката \\
\midrule
$s$ & $0$  & $(s,a)$, $(s,b)$          & $a{:}4(s)$, $b{:}2(s)$, $c{:}\infty$, $d{:}\infty$ \\
$b$ & $2$  & $(b,a)$, $(b,c)$, $(b,d)$ & $a{:}3(b)$, $c{:}10(b)$, $d{:}12(b)$ \\
$a$ & $3$  & $(a,c)$                   & $c{:}8(a)$, $d{:}12(b)$ \\
$c$ & $8$  & $(c,d)$                   & $d{:}10(c)$ \\
$d$ & $10$ & няма изходящи ребра       & $Q=\varnothing$ \\
\bottomrule
\end{tabular}
\end{center}

Върховете се извличат в реда $s,b,a,c,d$, тоест по нарастващи окончателни оценки $0,2,3,8,10$, а не по номер или по брой ребра. Три пъти се задейства \verb|Decrease-Key|: оценката на $a$ пада от $4(s)$ на $3(b)$, на $c$ --- от $10(b)$ на $8(a)$, и на $d$ --- от $12(b)$ на $10(c)$. Обърнете внимание, че $a$ е бил с оценка $4$, докато още е чакал в $Q$; именно затова окончателна става само оценката на извлечения връх.

Полученият резултат е
\[
\delta(s,s)=0,\quad
\delta(s,a)=3,\quad
\delta(s,b)=2,\quad
\delta(s,c)=8,\quad
\delta(s,d)=10,
\]
а указателите $\pi$ дават най-късия път до $d$:
\[
s\to b\to a\to c\to d.
\]
\end{example}

\begin{supp}[Бележка]
Алгоритъмът на Дийкстра не е коректен при отрицателни тегла. Причината е, че връх може да бъде окончателно избран, а по-късно да се открие по-лек път до него през ребро с отрицателно тегло.
\end{supp}

\begin{example}[Дийкстра греши при отрицателно тегло]
Нека графът има върхове $s,a,b$ и ребра
\[
(s,a)\text{ с тегло }2,\qquad
(s,b)\text{ с тегло }3,\qquad
(b,a)\text{ с тегло }-2.
\]
Отрицателен цикъл няма --- графът е ацикличен. Истинското разстояние е
\[
\delta(s,a)=3+(-2)=1.
\]
Алгоритъмът обаче извлича първо $s$ и получава $a{:}2$, $b{:}3$. След това минималната оценка в $Q$ е тази на $a$, затова $a$ се извлича и попада в $D$ с $a.d=2$. Когато по-късно се обработи $b$ и се релаксира $(b,a)$, проверката $y\in Q$ е лъжа, защото $a$ вече не е в опашката, и оценката остава $2\ne 1$.

Точно това не може да се случи при неотрицателни тегла: тогава удължаването на маршрут никога не го прави по-лек, което е и същината на доказателството по-долу.
\end{example}

\subsection{Коректност}\label{s:8.3.5}

\begin{keythm}[Коректност на Дийкстра]
Нека $G$ е ориентиран тегловен граф без отрицателни тегла и $s$ е негов връх. След завършване на алгоритъма на Дийкстра за всеки връх $v$ е изпълнено
\[
v.d=\delta(s,v).
\]
\end{keythm}

\begin{proof}
Ще докажем, че когато връх $v$ бъде добавен в $D$, вече е изпълнено
\[
v.d=\delta(s,v).
\]

Да допуснем противното и нека $v$ е първият връх, добавен в $D$ с неточна оценка. От свойството на оценките може да има само случая
\[
v.d>\delta(s,v).
\]
Връх $v$ не е $s$, защото $s.d=0=\delta(s,s)$.

Нека $p$ е най-къс път от $s$ до $v$. В момента непосредствено преди добавянето на $v$ в $D$ разглеждаме първия връх $u$ по $p$, който не принадлежи на $D$. Нека $a$ е предшественикът на $u$ по $p$. Тогава $a\in D$ и при обработването на $a$ реброто $(a,u)$ е било релаксирано.

Тъй като $a$ е добавен преди $v$, по избора на $v$ имаме
\[
a.d=\delta(s,a).
\]
От лемата за релаксация по най-къс път следва
\[
u.d=\delta(s,u).
\]

Върховете $u$ и $v$ са извън $D$, а $v$ е избран от приоритетната опашка с минимална оценка. Следователно
\[
v.d\leq u.d.
\]
Понеже теглата са неотрицателни и $u$ предхожда $v$ по най-късия път $p$, имаме
\[
\delta(s,u)\leq\delta(s,v).
\]
А от предположението $v.d>\delta(s,v)$ получаваме
\[
v.d\leq u.d=\delta(s,u)\leq\delta(s,v)<v.d,
\]
което е противоречие. Следователно всеки връх се добавя в $D$ с точна оценка. При завършване всички върхове са обработени.
\end{proof}

\subsection{Сложност}\label{s:8.3.6}

Нека $n=|V|$ и $m=|E|$. Времето зависи от реализацията на приоритетната опашка.

\begin{itemize}
\item При неупорядочен масив намирането и извличането на минимум струва $\Theta(n)$ на итерация. Общата сложност е
\[
\Theta(n^2+m),
\]
обикновено записвана като $\Theta(n^2)$.

\item При двоична пирамида извличането на минимум и намаляването на ключ струват $\Theta(\log n)$. Има най-много $n$ извличания и до $m$ релаксации, които могат да намалят ключ. Получава се
\[
\Theta((n+m)\log n).
\]

\item При по-сложни реализации на приоритетни опашки анализът може да бъде по-добър за операцията по намаляване на ключа, но основната идея на алгоритъма не се изменя.
\end{itemize}

\section{Най-къси пътища в тегловен DAG}\label{s:8.4}

\subsection{Алгоритъм}\label{s:8.4.1}

Ориентираният ацикличен граф, или DAG, допуска топологична наредба на върховете: за всяко ребро $(u,v)$ връх $u$ се намира преди $v$ в наредбата. Това позволява върховете да се обработят така, че преди даден връх да са обработени всички възможни предшественици по път от източника.

Алгоритъмът решава SSSP в тегловен DAG и допуска отрицателни тегла.

\begin{verbatim}
DAG-SSSP(G, w, s)
    Toposort(G)
    ISS(G, s)
    for each x in topological order
        for each y in adj(x)
            Relax((x, y))
\end{verbatim}

\subsection{Коректност}\label{s:8.4.2}

\begin{keythm}[Коректност на алгоритъма за DAG]
Нека $G$ е тегловен ориентиран ацикличен граф и $s\in V$. След завършването на \texttt{DAG-SSSP} за всеки връх $v$ е изпълнено
\[
v.d=\delta(s,v).
\]
\end{keythm}

\begin{proof}
Нека върховете са разгледани в топологична наредба. За всеки връх, стоящ преди $s$, няма път от $s$ до него, защото всеки път следва топологичната наредба. Следователно за него
\[
v.d=+\infty=\delta(s,v).
\]
За $s$ след инициализацията е изпълнено
\[
s.d=0=\delta(s,s),
\]
тъй като в DAG няма цикли, а следователно няма и отрицателни цикли.

Нека $v$ е произволен връх след $s$ в топологичната наредба. Всеки вътрешен връх на път от $s$ до $v$ стои преди $v$. Нека
\[
U=\{u\in V\mid (u,v)\in E \text{ и има път от }s\text{ до }u\}.
\]
До момента, в който започне обработването на $v$, всички върхове от $U$ вече са обработени. По индуктивното предположение за тях е вярно
\[
u.d=\delta(s,u).
\]
При релаксиране на всички входящи към $v$ възможности, реализирани чрез обработката на техните начала, оценката става
\[
v.d=\min_{u\in U}\{\delta(s,u)+\omega((u,v))\}.
\]
По оптималната подструктура това е точно $\delta(s,v)$. Ако $U=\emptyset$, връх $v$ не е достижим от $s$; тогава и двете стойности са $+\infty$.
\end{proof}

\subsection{Сложност и предимства}\label{s:8.4.3}

Топологичното сортиране има сложност $\Theta(n+m)$. Вложеният цикъл обхожда всеки връх и всяко ребро най-много веднъж, затова общата сложност е
\[
\Theta(n+m).
\]

Алгоритъмът има две важни предимства пред алгоритъма на Дийкстра върху DAG:

\begin{itemize}
\item работи с отрицателни тегла, защото в DAG няма цикли и следователно няма отрицателни цикли;
\item има линейна сложност $\Theta(n+m)$, която е асимптотично по-добра от сложността на Дийкстра с приоритетна опашка.
\end{itemize}

\begin{remark}
Чрез обръщане на неравенството в операцията за релаксация алгоритъмът за DAG може аналогично да намира най-дълги пътища. Това е възможно именно поради липсата на цикли.
\end{remark}

\section{Алгоритъм на Белман--Форд}\label{s:8.5}

\subsection{Вариант на задачата и псевдокод}\label{s:8.5.1}

Алгоритъмът на Белман--Форд решава SSSP в ориентиран тегловен граф, който може да съдържа отрицателни ребра. Освен това открива наличие на отрицателен цикъл, достижим от началния връх.

\begin{verbatim}
Bellman-Ford(G, w, s)
    ISS(G, s)
    for i <- 1 to n - 1
        for each (x, y) in E(G)
            Relax((x, y))
    for each (x, y) in E(G)
        if y.d > x.d + w((x, y))
            return False
    return True
\end{verbatim}

Ако алгоритъмът върне \texttt{True}, отрицателен цикъл, достижим от $s$, не е открит и оценките представляват най-късите разстояния. Ако върне \texttt{False}, съществува достижим отрицателен цикъл.

\subsection{Коректност при липса на достижими отрицателни цикли}\label{s:8.5.2}

\begin{lem}
Нека в графа няма отрицателен цикъл, достижим от $s$. След изпълнение на първите $n-1$ обхождания на всички ребра от алгоритъма на Белман--Форд за всеки връх $v$ е изпълнено
\[
v.d=\delta(s,v).
\]
\end{lem}

\begin{proof}
При липса на отрицателни цикли, достижими от $s$, могат да се изберат прости най-къси пътища от $s$ до достижимите върхове и техните предшественици да се разглеждат като дърво на най-късите пътища с корен $s$. Недостижимите върхове запазват оценка $+\infty$: ако към такъв връх имаше ребро от достижим връх, той също би бил достижим.

Ще използваме индукция по нивата на това дърво. Преди първото обхождане е вярно
\[
s.d=0=\delta(s,s).
\]
Да допуснем, че след някакъв брой обхождания всички върхове до ниво $i-1$ имат точни оценки. Нека $t$ е връх на ниво $i$, а $v$ е неговият родител в дървото на най-късите пътища. При следващото обхождане реброто $(v,t)$ се релаксира. Тъй като
\[
v.d=\delta(s,v),
\]
от лемата за релаксация по най-къс път следва
\[
t.d=\delta(s,t).
\]
Това равенство се запазва.

Всеки прост път съдържа най-много $n-1$ ребра. Следователно след $n-1$ пълни обхождания са достигнати всички нива на дървото на най-късите пътища и всички оценки са точни.
\end{proof}

\begin{lem}
Ако в графа няма отрицателен цикъл, достижим от $s$, то след първите $n-1$ обхождания за всяко ребро $(x,y)$ е изпълнено
\[
y.d\leq x.d+\omega((x,y)).
\]
\end{lem}

\begin{proof}
Ако $x$ е достижим от $s$, тогава и $y$ е достижим и след първите $n-1$ обхождания имаме $x.d=\delta(s,x)$ и $y.d=\delta(s,y)$. Твърдението следва от неравенството за ребро. Ако $x$ не е достижим, тогава $x.d=+\infty$ и неравенството е тривиално.
\end{proof}

\begin{thm}
Ако в графа няма отрицателен цикъл, достижим от $s$, алгоритъмът на Белман--Форд връща \texttt{True} и за всеки връх $v$ е изпълнено
\[
v.d=\delta(s,v).
\]
\end{thm}

\begin{proof}
Точността на оценките след първия двоен цикъл следва от първата лема в този подраздел. От следващата лема тогава за никое ребро условието в последния цикъл не е изпълнено, поради което алгоритъмът достига \texttt{return True}.
\end{proof}

\subsection{Откриване на отрицателен цикъл}\label{s:8.5.3}

\begin{lem}
Ако в графа има отрицателен цикъл, достижим от $s$, то при проверката след $n-1$ обхождания съществува ребро $(x,y)$, за което
\[
y.d>x.d+\omega((x,y)).
\]
\end{lem}

\begin{proof}
Да допуснем противното: за всяко ребро $(x,y)$ е изпълнено
\[
y.d\leq x.d+\omega((x,y)).
\]
Нека отрицателният цикъл е
\[
v_1,e_1,v_2,e_2,\ldots,v_k,e_k,v_1.
\]
Тъй като цикълът е достижим от $s$, след извършените обхождания всички негови върхове имат крайни оценки. Прилагаме предположеното неравенство за всички ребра на цикъла:
\[
d(v_2)\leq d(v_1)+\omega(e_1),
\]
\[
d(v_3)\leq d(v_2)+\omega(e_2),
\]
\[
\cdots
\]
\[
d(v_1)\leq d(v_k)+\omega(e_k).
\]
Като съберем неравенствата, всички оценки се съкращават и получаваме
\[
0\leq \omega(e_1)+\omega(e_2)+\cdots+\omega(e_k).
\]
Това противоречи на отрицателното тегло на цикъла.
\end{proof}

\begin{keythm}[Коректност на Белман--Форд]
Алгоритъмът на Белман--Форд връща \texttt{True} точно когато в графа няма отрицателен цикъл, достижим от $s$; в този случай за всеки връх $v$ е изпълнено
\[
v.d=\delta(s,v).
\]
Ако върне \texttt{False}, има достижим отрицателен цикъл.
\end{keythm}

\begin{proof}
Ако няма достижим отрицателен цикъл, всеки достижим връх има най-къс прост път с най-много $|V|-1$ ребра. По индукция по дължината на такъв път след $k$ пълни обхода на ребрата е получена правилната стойност за всеки връх с най-къс път от най-много $k$ ребра: поредният обход релаксира последното ребро след вече получения префикс. Оценките не могат да са по-малки от истинското разстояние, защото всяка крайна оценка е тегло на реален маршрут. След $|V|-1$ обхода всички оценки са точни и никое ребро не допуска подобрение; алгоритъмът връща \texttt{True}.

Ако има достижим отрицателен цикъл, всички негови върхове имат крайни оценки след тези обходи. Ако никое негово ребро не допуска подобрение, за всяко $(u,v)$ от цикъла е $v.d\le u.d+w(u,v)$. Сумирането съкращава оценките и дава $0\le w(C)$, противоречие. Значи проверката открива подобрение и връща \texttt{False}. Недостижимите върхове остават с $+\infty$ и не предизвикват такова подобрение.
\end{proof}

\subsection{Сложност}\label{s:8.5.4}

Външният цикъл се изпълнява $n-1$ пъти, а при всяко изпълнение се обхождат всички $m$ ребра. Последната проверка има сложност $\Theta(m)$. Следователно времевата сложност е
\[
\Theta(nm).
\]
В най-лошия случай, когато $m=\Theta(n^2)$, тя е кубична:
\[
\Theta(n^3).
\]

\section{Алгоритъм на Флойд--Уоршал}\label{s:8.6}

\subsection{Вариант на задачата}\label{s:8.6.1}

Алгоритъмът на Флойд--Уоршал решава задачата APSP: намира най-късите разстояния между всички двойки върхове. Графът е ориентиран и е номериран с върхове
\[
V=\{1,\ldots,n\}.
\]

Входът се представя с матрица на теглата $\omega$, където
\[
\omega[i,j]=
\begin{cases}
0, & i=j,\\
\omega((i,j)), & i\ne j \text{ и }(i,j)\in E,\\
+\infty, & i\ne j \text{ и }(i,j)\notin E.
\end{cases}
\]

\subsection{Динамично програмиране}\label{s:8.6.2}

За път $p$ нека $\operatorname{intern}(p)$ е множеството от неговите вътрешни върхове. Нека $D^{(k)}[i,j]$ означава теглото на най-къс път от $i$ до $j$, чиито вътрешни върхове са само измежду
\[
\{1,\ldots,k\}.
\]

При $k=0$ не се разрешават вътрешни върхове, затова
\[
D^{(0)}=\omega.
\]

За $k>0$ има две възможности за най-къс път от $i$ до $j$ с разрешени вътрешни върхове от $\{1,\ldots,k\}$:

\begin{itemize}
\item връх $k$ не участва като вътрешен връх; тогава стойността е $D^{(k-1)}[i,j]$;
\item връх $k$ участва като вътрешен връх; тогава пътят се разделя на път от $i$ до $k$ и път от $k$ до $j$, чиито вътрешни върхове са от $\{1,\ldots,k-1\}$.
\end{itemize}

Следователно рекурентната формула е
\[
D^{(k)}[i,j]
=
\min\left\{
D^{(k-1)}[i,j],
D^{(k-1)}[i,k]+D^{(k-1)}[k,j]
\right\}.
\]

\subsection{Псевдокод}\label{s:8.6.3}

\begin{verbatim}
Floyd-Warshall(w)
    D <- w
    for k <- 1 to n
        for i <- 1 to n
            for j <- 1 to n
                D[i, j] <- min(D[i, j], D[i, k] + D[k, j])
    return D
\end{verbatim}

\subsection{Коректност}\label{s:8.6.4}

\begin{keythm}[Коректност на Флойд--Уоршал]
При отсъствие на отрицателни цикли алгоритъмът на Флойд--Уоршал връща матрица $D$, за която
\[
D[i,j]=\delta(i,j)
\]
за всички $i,j\in\{1,\ldots,n\}$.
\end{keythm}

\begin{proof}
Ще докажем с индукция по $k$, че след завършване на итерацията за $k$ матрицата съдържа точно $D^{(k)}$.

За $k=0$ това е вярно по дефиницията на входната матрица $\omega$: позволени са само пътища без вътрешни върхове.

Нека твърдението е вярно за $k-1$. За всяка двойка $i,j$ рекурентната формула разглежда точно двата възможни случая: връх $k$ не помага или $k$ е вътрешен връх на пътя. Във втория случай оптималният път се разделя на два подпътя с вътрешни върхове, номерирани най-много с $k-1$. Следователно присвояването изчислява точно $D^{(k)}[i,j]$.

При $k=n$ всички върхове са разрешени като вътрешни. При липса на отрицателни цикли всеки най-къс път може да бъде избран прост, така че получената стойност е $\delta(i,j)$.
\end{proof}

\subsection{Сложност по време и памет}\label{s:8.6.5}

Трите вложени цикъла извършват $\Theta(n^3)$ операции. Следователно времевата сложност е
\[
\Theta(n^3).
\]

Буквалното съхраняване на всички матрици
\[
D^{(0)},D^{(1)},\ldots,D^{(n)}
\]
изисква
\[
\Theta(n^3)
\]
памет. Достатъчно е обаче да се пазят само две работни матрици: една за $D^{(k-1)}$ и една за $D^{(k)}$. Така използваната памет става
\[
\Theta(n^2).
\]

Възможна е и реализация на място, при която се използва една матрица $D$ и се изпълнява
\[
D[i,j]\leftarrow
\min\{D[i,j],D[i,k]+D[k,j]\}.
\]
Тогава извън входната матрица е нужна $\Theta(1)$ допълнителна памет, но входната матрица се презаписва. При липса на отрицателни цикли случаите $i=k$ или $j=k$ не пречат на коректността, защото диагоналните стойности са нула и връх $k$ не може да бъде вътрешен връх на прост път, който започва или завършва в $k$.

\section{Изпитен фокус}\label{s:8.7}

\begin{itemize}
\item Да се дефинират тегло на път, $\delta(u,v)$, отрицателен цикъл и четирите варианта на задачата за най-къси пътища.

\item Да се обясни защо отрицателен цикъл, разположен по път от $u$ до $v$, прави стойността $\delta(u,v)$ равна на $-\infty$, както и защо при липса на достижим отрицателен цикъл може да се избере прост най-къс път.

\item Да се възпроизведат и използват оптималната подструктура на най-късите пътища, инициализацията \texttt{ISS}, операцията \texttt{Relax} и инвариантът $v.d\geq\delta(s,v)$.

\item За Дийкстра: да се посочат условието за неотрицателни тегла, псевдокодът, множеството $D$, операциите \verb|Extract-Min| и \verb|Decrease-Key| на приоритетната опашка и ролята на проверката $y\in Q$, доказателствената идея с първия неправилно избран връх и сложностите при различни реализации на опашката.

\item Да се посочи контрапример, показващ, че при отрицателно тегло Дийкстра дава грешен резултат дори без отрицателен цикъл.

\item За DAG-SSSP: да се посочат топологичната наредба, релаксирането на всички ребра в тази наредба, доказателството по реда на върховете, сложността $\Theta(n+m)$ и предимствата пред Дийкстра.

\item За Белман--Форд: да се възпроизведат $n-1$ пълни обхождания на всички ребра, финалната проверка, доказателството чрез нивата в дърво на най-късите пътища и доказателството за откриване на отрицателен цикъл чрез сумиране на неравенства по цикъла.

\item За Флойд--Уоршал: да се дефинира $D^{(k)}[i,j]$, да се изведе рекурентната формула според това дали връх $k$ участва, да се даде псевдокодът, индуктивното доказателство и сложностите $\Theta(n^3)$ по време и $\Theta(n^2)$ памет при реализация с една или две матрици.
\end{itemize}
